\section*{Chapitre \ref{ch:b3:pericyclic} --- Réactions péricycliques}

\begin{solution}{exo:b3:pericyclic:1}
Diels--Alder~: \omterm{def:b3:pericyclic:families}{cycloaddition} [4+2]. Ouverture du cyclobutène~:
électrocyclique. Cope~: sigmatropique [3,3]. Perte de $\ce{SO2}$ par le sulfolène~: chélétropique. Ozone et alcène~: \omterm{def:b3:pericyclic:dipolar}{cycloaddition
1,3-dipolaire} ([3+2]). Migration $[1,5]$-H~: sigmatropique.
\end{solution}

\begin{solution}{exo:b3:pericyclic:2}
Hexatriène (6 électrons)~: thermique \omterm{def:b3:pericyclic:rotation-modes}{disrotatoire}, photochimique
\omterm{def:b3:pericyclic:rotation-modes}{conrotatoire}. Butadiène (4)~: thermique \omterm{def:b3:pericyclic:rotation-modes}{conrotatoire}, photochimique
\omterm{def:b3:pericyclic:rotation-modes}{disrotatoire}.
\end{solution}

\begin{solution}{exo:b3:pericyclic:3}
Quatre électrons, thermique~: ouverture \omterm{def:b3:pericyclic:rotation-modes}{conrotatoire}. Les deux groupes
méthyle, sur des faces opposées du cycle, tournent tous deux vers
l'extérieur~: (\textit{E},\textit{E})-hexa-2,4-diène.
\end{solution}

\begin{solution}{exo:b3:pericyclic:4}
La réaction [2+2] \omterm{def:b3:pericyclic:faciality}{suprafaciale} a $4q$ électrons~: dans l'état fondamental, la HOMO d'un
éthène et la LUMO de l'autre se recouvrent avec une extrémité liante et une
extrémité antiliante. Dans l'état excité, la $\pi^*$ simplement occupée d'une molécule joue le rôle de la
HOMO, et les deux extrémités s'accordent.
\end{solution}

\begin{solution}{exo:b3:pericyclic:5}
Le plan miroir est conservé. Hexatriène~: $\psi_1$ S, $\psi_2$ A, $\psi_3$ S occupées~; $\psi_4$ A, $\psi_5$ S, $\psi_6$ A vides.
Cyclohexadiène, dans l'ordre~: $\sigma$ S, $\pi_1$ S, $\pi_2$ A occupées~; $\pi_3^*$ S, $\pi_4^*$ A, $\sigma^*$ A vides. En reliant par
symétrie dans l'ordre, $\psi_1 \to \sigma$, $\psi_2 \to \pi_2$, $\psi_3 \to \pi_1$~: des occupées vers des occupées, la fermeture \omterm{def:b3:pericyclic:rotation-modes}{disrotatoire}
est thermiquement permise.
\end{solution}

\begin{solution}{exo:b3:pericyclic:6}
Six électrons sous irradiation~: \omterm{def:b3:pericyclic:rotation-modes}{conrotatoire}, si bien que les deux groupes
méthyle terminaux aboutissent sur des faces opposées~: \textit{trans}-5,6-diméthylcyclohexa-1,3-diène (la réaction thermique donne
l'isomère \textit{cis}).
\end{solution}

\begin{solution}{exo:b3:pericyclic:7}
L'hydrogène porté par C5 migre vers C1 du diène, son voisin sur le cycle,
via un état de transition à six électrons ($\sigma2_s + \pi4_s$), de façon \omterm{def:b3:pericyclic:faciality}{suprafaciale}~: une migration
donne le 1-méthyl-, une seconde le 2-méthylcyclopentadiène. Une chaîne de
cinq carbones permet à l'hydrogène de rester sur une même face, ce
qu'impose de toute façon le petit cycle.
\end{solution}

\begin{solution}{exo:b3:pericyclic:8}
$\pi2_s$ (alcène) $+ \pi2_a$ (C=C du cétène, attaquée sur des faces opposées grâce à
l'orbitale $\pi^*$ C=O perpendiculaire)~: le nombre de composantes $(4q + 2)_s$ vaut un (l'alcène), impair~:
thermiquement permise.
\end{solution}

\begin{solution}{exo:b3:pericyclic:9}
Dans la chaise, la C=C du groupe vinyle (atomes 1--2), l'oxygène (3), le carbone $\ce{OCH2}$
(4) et la C=C allylique (5--6)~; la liaison O--C se rompt et la liaison C1--C6
se forme~: le produit est le pent-4-énal,
$\ce{OHC-CH2-CH2-CH=CH2}$.
\end{solution}

\begin{solution}{exo:b3:pericyclic:10}
Huit électrons, c'est $4q$~; avec une composante \omterm{def:b3:pericyclic:faciality}{antarafaciale},
l'arrangement est de topologie de Möbius, aromatique avec $4q$ électrons~: l'état de transition est aromatique et la
réaction thermiquement permise (comme la fermeture \omterm{def:b3:pericyclic:rotation-modes}{conrotatoire} de
l'octatétraène).
\end{solution}

\begin{solution}{exo:b3:pericyclic:11}
En associant les plus grands coefficients, on relie N3 au carbone
terminal~: le triazole 1,4-disubstitué. Sans catalyseur, les coefficients
diffèrent peu et les deux isomères 1,4 et 1,5 se forment~; le cuivre(I)
rend la réaction séquentielle, via un acétylure de cuivre, et donne
l'isomère 1,4 seul.
\end{solution}

\begin{solution}{exo:b3:pericyclic:12}
L'hydrogène fait le pont entre les extrémités de la SOMO d'un radical de $j$ atomes, $k = (j + 1)/2$. Ses coefficients
terminaux sont proportionnels à $\sin(k\pi/(j + 1))$ et $(-1)^{k+1}\sin(k\pi/(j + 1))$. Pour $j = 7$, $k = 4$~: signes opposés, si bien que
les lobes de même signe sont sur des faces opposées aux deux extrémités~: \omterm{def:b3:pericyclic:faciality}{antarafacial}. Pour
$j = 5$, $k = 3$~: même signe sur la même face~: \omterm{def:b3:pericyclic:faciality}{suprafacial}.
\end{solution}

\begin{solution}{pb:b3:pericyclic:1}
\textbf{1.} La liaison $\sigma$ C9--C10 du cycle B~; avec le diène 5,7, elle donne un hexatriène,
C10=C5--C6=C7--C8=C9.

\textbf{2.} Six électrons (les deux liaisons $\pi$ et la liaison $\sigma$)~: une ouverture de cycle électrocyclique.

\textbf{3.} Par voie thermique, un cyclohexadiène et l'hexatriène ouvert
correspondant ne s'équilibrent que lentement, à travers une barrière
élevée, et le cycle est le côté le plus stable~; le diène absorbe la
lumière ultraviolette, et son état excité s'ouvre en quelques
picosecondes.

\textbf{4.} \omterm{def:b3:pericyclic:rotation-modes}{Conrotatoire}~: dans l'état excité, la LUMO simplement occupée du système triène
($k = 4$) a des extrémités de signes opposés.

\textbf{5.} La liaison C6=C7 provient du cycle~: ses deux prolongements de
chaîne, C5 et C8, étaient du même côté de cette liaison, et y restent~:
\textit{Z}.

\textbf{6.} L'ouverture thermique, permise seulement de façon
\omterm{def:b3:pericyclic:rotation-modes}{disrotatoire}, a une barrière bien au-delà de ce que peut fournir la chaleur
du corps, et conduirait vers le haut, vers le triène, moins stable.

\textbf{7.} De C19 (le méthyle porté par C10) vers C9~: une migration sigmatropique $[1,7]$-H.

\textbf{8.} Huit~: les six électrons $\pi$ du triène et les deux de la liaison C--H.

\textbf{9.} $\sigma2 + \pi6$. Si tout était \omterm{def:b3:pericyclic:faciality}{suprafacial}, les deux seraient s, chacune comptant comme une composante $(4q + 2)_s$~:
nombre pair, interdit. Avec le triène \omterm{def:b3:pericyclic:faciality}{antarafacial}, $\sigma2_s + \pi6_a$~: une seule composante comptée, nombre impair,
permis.

\textbf{10.} Sept atomes dans une chaîne hélicoïdale de configuration \textit{Z} permettent à l'hydrogène de passer d'une face à
l'autre~; une chaîne de trois atomes ne peut pas se tordre autant.

\textbf{11.} Möbius (une composante \omterm{def:b3:pericyclic:faciality}{antarafaciale}), avec huit électrons, $4q$~: aromatique.

\textbf{12.} Elle est thermiquement permise~; sa barrière est franchie lentement à la température du corps.

\textbf{13.} Six électrons sous irradiation~: de nouveau \omterm{def:b3:pericyclic:rotation-modes}{conrotatoire}, mais
la fermeture peut tourner les extrémités dans l'autre sens et placer H9 et
le méthyle C19 sur les autres faces~: un stéréoisomère à cycle fermé, le
lumistérol.

\textbf{14.} La liaison centrale étant \textit{E}, C19 et C9 se trouvent de part et d'autre de la chaîne et loin
l'un de l'autre~: l'hydrogène ne peut pas l'atteindre.

\textbf{15.} La prévitamine D$_3$ elle-même absorbe et se convertit en lumistérol et en tachystérol,
qui ne donnent pas de vitamine D~; un mélange photostationnaire est atteint,
ce qui plafonne la production.

\textbf{16.} L'ouverture a lieu dans l'état excité, en quelques
femtosecondes à picosecondes~; la migration est une réaction thermique avec une barrière de \qty{85}{kJ/mol}.

\textbf{17.} $t_{1/2} = \ln 2/k = 0.693/\qty{0.023}{h^{-1}} = \qty{30}{h}$.

\textbf{18.} $1 - \eu^{-0.023 \times 24} = 0.42$.

\textbf{19.} $\ln 10/k = 2.303/\qty{0.023}{h^{-1}} = \qty{100}{h}$.

\textbf{20.} $k(\qty{293.15}{K}) = k(\qty{310.15}{K})\,\eu^{-(E_a/R)(1/293.15 - 1/310.15)} = 0.023 \times \eu^{-1.91} = \qty{3.4e-3}{h^{-1}}$.

\textbf{21.} $2.303/\qty{3.4e-3}{h^{-1}} = \qty{680}{h}$, soit environ 28 jours.

\textbf{22.} L'étape thermique est sept fois plus rapide à \qty{37}{\celsius} qu'à \qty{20}{\celsius}, si bien que la
prévitamine formée en une matinée au soleil est convertie en quelques
jours.

\textbf{23.} Le squelette stéroïde, avec ses nombreux stéréocentres, est
fourni tout fait par le stérol~; la lumière et la chaleur réalisent ensuite
les deux étapes péricycliques exactement comme dans la peau.

\textbf{24.} À la température du corps, 90\,\% de la prévitamine D$_3$ deviennent de la vitamine D$_3$ en environ
\qty{100}{h}, soit quatre jours.
\end{solution}
